\section{探索激励与长度控制：质量-成本联合优化}

讲者反复讨论探索（exploration）与成本控制之间的张力。模型若被鼓励过度探索，会出现超长轨迹与计算爆炸；若惩罚过强，又会抑制有效搜索，导致质量提升停滞。

这可以形式化为一个联合目标：
\[
\max_{\pi} \ \mathbb{E}_{\tau \sim \pi}\left[R(\tau)\right]
- \lambda \cdot \mathbb{E}[C_{\text{token}}(\tau)]
- \mu \cdot \mathbb{E}[T_{\text{latency}}(\tau)]
- \beta \cdot D_{\mathrm{KL}}(\pi \Vert \pi_{\text{ref}})
\]
其中成本项和时延项在产品场景中与质量同等重要。

\begin{importantbox}{课堂中的关键工程经验}
评估时只报告 pass rate 不足以比较模型。必须同时报告 rollout length / test-time compute，否则高分可能只是因为给了更长推理预算。
\end{importantbox}

字幕中有一段非常直接的观点：如果一个方案一小时完成，另一个方案一分钟完成且质量接近，用户一定更偏好后者。这提醒我们：agent 模型不是离线竞赛模型，必须在可接受时延内稳定交付。

\begin{knowledgebox}{长度控制常见做法}
\begin{itemize}
    \item 训练时引入 token cost penalty；
    \item 在奖励中加入 ``短解优先'' 偏好；
    \item 对超长 rollout 设 hard stop 与回退机制；
    \item 评估时统一预算，做 apples-to-apples 对比。
\end{itemize}
\end{knowledgebox}

\begin{warningbox}{过度 length penalty 的副作用}
若惩罚过高，模型会避免必要探索，出现 ``短而错'' 的策略塌缩。正确做法是分阶段调节：前期鼓励探索，中后期强化效率。
\end{warningbox}

\subsection{本章小结}
Agent 训练的目标不是 ``无限算力下最优''，而是 ``给定预算下最优''。质量与成本必须被联合建模并同时评估。

